\chapter{ネコをどう見分けるか}
% full version: chapters/12_recognition.tex

\pencilfig{12}{\pencilart{12}}{大きなネコと小さなネコは、目のセンサ上ではまったく別の絵なのに、答えは1つ。}

絵の中のネコは、左にいることも右にいることも、近くにいることも遠くにいることも、日なたにいることも日陰にいることもある。目のセンサにはそのたびにまったく違う模様が落ちるのに、あなたは一瞬でネコだとわかる。これが認識の難しさである。同じものを、何千通りもの見え方の中から見分けなければならない。

\section*{組立ライン}

脳はこの問題を、およそ10段からなる組立ラインで解く。信号はそれを0.15秒で通り抜ける。各段では2つのことが行われる。まず「AND」。必要な部品が必要な並び方をしていれば、たとえば2本の線分が角をなして交われば、ニューロンが鳴る。次に「すべての場所でのOR」。近くのどこかに角があれば、正確な場所に関係なく次のニューロンが鳴る。1つ目の演算は認識を選り好みさせ、2つ目はずれに寛容にする。段を重ねるごとに部品は部分に、部分は物体にまとまり、ずれや大きさの違いは見逃される。

\pencilfig{112}{%
\foreach \i/\y/\cx/\cy in {1/1.2/0.15/0.3,2/0.3/0.3/0.1,3/-0.6/0.05/0.05}{
  \draw[penlite] (0,\y-0.3) rectangle (0.6,\y+0.3);
  \draw[pcGray, line width=0.8pt] (\cx,\y-0.3+\cy+0.35) -- (\cx,\y-0.3+\cy) -- (\cx+0.3,\y-0.3+\cy);
  \draw[pen=pcBlue, wash=pcBlue] (1.9,\y) circle (0.2);
  \draw[arr] (0.65,\y) -- (1.65,\y);
  \draw[arr=pcBlue] (2.12,\y) -- (3.62,{0.3+0.12*(2-\i)});}
\draw[pen=pcBlue, wash=pcBlue] (3.9,0.3) circle (0.25);
\draw[arr] (4.2,0.3) -- (5.75,0.3); \node[lbl, anchor=south] at (4.97,0.35) {さらに段};
% a small cat
\draw[penlite] (6.3,0.3) circle (0.32);
\draw[penlite] (6.03,0.5) -- (6.07,0.82) -- (6.23,0.6); \draw[penlite] (6.57,0.5) -- (6.53,0.82) -- (6.37,0.6);
\node[lbl, anchor=north] at (0.3,-1.0) {絵}; \node[lbl, anchor=north, align=center] at (1.9,-1.0) {AND：角が\\この場所に};
\node[lbl, anchor=north, align=center] at (3.9,-1.0) {OR：角が\\どこかに}; \node[lbl, anchor=north] at (6.3,-1.0) {ネコ};
}{組立ラインの1段：「AND」は1か所で角を組み立て、「OR」はそれがどこにあるかを問わない。}

スマートフォンで画像を認識する畳み込みニューラルネットワークは、まさにこの仕組みを写したものだ。そしてお返しもしてくれる。物体の認識を学習したネットワークは、視覚の上位段のニューロンがどう応答するかをかなりよく予測する。

\section*{動画から学ぶ}

だが、誰が脳に教えたのか。「ネコ」とラベルのついた絵を100万枚見せてくれた人はいない。有力な答えは動画である。世界は途切れなく流れ、一瞬のうちに物体はほとんど変わらず、変わるのは位置や照明だ。いま見えているものと一瞬前に見えていたものの間の結合が強まるなら、ネットワークは同じ物体の異なる見え方を自分で結びつける。ゆっくり変わるものが物体で、速く変わるものが状況だ。眼球運動の最中に物体をこっそりすり替える実験があり、するとニューロンは別々の物体を同じものとみなし始めた。一般的な規則としては、これはまだ仮説である。

\section*{錯覚は仕組みの指紋}

錯覚は故障ではなく、マシンの作りの痕跡である。滝を長く見つめてから動かない岩を見ると、岩が上へ「はい上がる」。「下」と「上」のチャネルの組が、片方の疲労でバランスを崩したのだ。ある人には白と金に、別の人には青と黒に見える有名なドレスは、その人の脳が暗黙のうちにどんな照明を仮定しているかで人を二分する。描かれた「回る蛇」が動いて見えるのは、明るい部分が暗い部分よりわずかに速く処理され、その時間差が動きとして読まれるからだ。そして飛び出して見えたりへこんで見えたりする立方体は、互いに競い合う2つのニューロン群である。あるモデルによれば、結果の入れ替わりは主にノイズが決めている。

興味深いことに、動画の次のフレームを予測するよう訓練したニューラルネットワークにも、蛇が回って「見える」。つまり、錯覚のいくつかは予測能力の避けられない代償なのだ。

\section*{脳対ニューラルネットワーク}

似てはいるが、完全に同じではない。脳は少数の例で足り、ラベルがほとんどなくても学び、すべてを20ワットでこなす。ニューラルネットワークは、人間には気づかないほどの画像の細工でだませる。ただ、こうした細工を一瞬だけ見せると、人間も少し惑わされる。この2つのマシンの本当の境界がどこにあるのかは、まだ調べられている最中である。

\fullversion{12}
